\section*{Appendix A: Retained Rubric Definition}

Two senior professional evaluators rated each retained segment on LQ, EXP, and
professional promptness (LAT in the archived rubric) using the same 0--3 integer-anchor rubric. Higher
promptness means the interpretation was experienced as \emph{less} delayed. The
aggregate target is their equal-weight mean, which produces .5-point increments.
Both evaluators rated the full cohort independently after shared written instructions,
example discussion, and a pilot calibration exercise; neither served as adjudicator.
The scoring platform enforced LQ$=0\Rightarrow$EXP$=0$ and promptness$=0$; such
lower-bound triples therefore reflect a rubric constraint rather than independent
dimension-level decisions.
Table~\ref{tab:retained-grid} reproduces the retained anchor criteria.

\begin{table}[h]
\centering
\scriptsize
\setlength{\tabcolsep}{2pt}
\caption{Retained segment-level rubric definition. Individual evaluators use the integer anchors 0--3; .5-point values occur only in the two-rater mean target.}
\label{tab:retained-grid}
\begin{tabular}{p{.25\columnwidth}p{.66\columnwidth}}
\toprule
Dimension & Criteria summary \\
\midrule
LQ & 3: meaning is preserved with no major omission or distortion. 2: some omissions or inaccuracies, but the main message remains intact. 1: major omissions or distortions, including misleading or reversed meaning. 0: meaning is largely absent, misleading, or reversed. \\
EXP & 3: fluent, idiomatic; minor errors only. 2: generally fluent with noticeable awkwardness/errors. 1: frequent disfluency/grammar issues affecting delivery. 0: breakdown-level disfluency impedes comprehension. \\
Promptness\newline (archived label: LAT) & 3: low perceived lag; well synchronized. 2: occasional lag; manageable. 1: frequent noticeable lag; disjointed flow. 0: severe persistent lag; difficult to follow. \\
\bottomrule
\end{tabular}
\setlength{\tabcolsep}{6pt}
\end{table}

\section*{Appendix B: Additional Audit Results}

\paragraph{Platform-rule audit.} Among the $2\times\CohortN=1{,}244$
individual evaluator--segment score triples in the retained cohort, 17 have
LQ$=0$. All 17 also have EXP$=0$ and promptness$=0$, with no violation of the
platform rule; all 17 were assigned by R05. Thus, 17 of the \CohortN\ aggregate
segments (2.7\%) contain one platform-constrained lower-bound triple.

\paragraph{Rater agreement.} On all \CohortN\ jointly rated segments, R05 assigned scores
$\{0,1,2,3\}$ with counts $(\RZeroCount,\ROneCount,\RTwoCount,\RThreeCount)$ and R06 with counts
$(\SZeroCount,\SOneCount,\STwoCount,\SThreeCount)$. The rows of Table~\ref{tab:rater-confusion} are R05 scores
and the columns are R06 scores. Exact agreement is \InterraterExactPercent\%; within-one-point
agreement is \InterraterWithinOnePercent\%; Pearson and Spearman correlations are \InterraterPearson\ and \InterraterSpearman;
ICC$(2,1)$ and quadratic-weighted kappa are \InterraterICC\ and \InterraterKappa.

% Generated by build_eacl_paper_results.py. Do not edit by hand.
\begin{table}[h]
\centering
\small
\caption{Promptness-score confusion matrix: R05 rows, R06 columns.}
\label{tab:rater-confusion}
\begin{tabular}{c|rrrr}
\toprule
R05 $\backslash$ R06 & 0 & 1 & 2 & 3 \\
\midrule
0 & 0 & 4 & 8 & 7 \\
1 & 0 & 3 & 41 & 28 \\
2 & 0 & 16 & 60 & 110 \\
3 & 0 & 10 & 78 & 257 \\
\bottomrule
\end{tabular}
\end{table}


\paragraph{Individual-rater human-label sensitivity.} Table~\ref{tab:rater-human}
reports selected held-out speech-group results using human LQ/EXP labels. These
are construct diagnostics rather than deployable automatic systems. R05's
same-rater reference is much higher than R06's, and each rater is differently
related to the onset signal, reinforcing the evaluator-calibration framing.

% Generated by build_eacl_paper_results.py. Do not edit by hand.
\begin{table*}[t]
\centering
\scriptsize
\caption{Individual-rater human-label audit. Models are fitted within the same outer speech-group folds; intervals are speech-group clustered 95\% CIs.}
\label{tab:rater-human}
\begin{tabular}{llccc}
\toprule
Target & System & Pearson $r$ & Spearman $r_S$ & MSE \\
\midrule
R05 & Delay & .106 [.014,.220] & .098 & .656 \\
R05 & Same-rater LQ+EXP+delay & .909 [.855,.944] & .836 & .113 \\
R05 & Structure+delay & .613 [.457,.691] & .546 & .406 \\
R05 & Full human reference & .911 [.859,.944] & .840 & .111 \\
R06 & Delay & .400 [.293,.488] & .403 & .293 \\
R06 & Same-rater LQ+EXP+delay & .552 [.496,.607] & .547 & .242 \\
R06 & Structure+delay & .442 [.373,.498] & .479 & .281 \\
R06 & Full human reference & .545 [.481,.610] & .552 & .245 \\
\bottomrule
\end{tabular}
\end{table*}


\paragraph{Cross-rater associations.} Table~\ref{tab:cross-rater} gives the
same-rater and cross-rater results summarized in Section~5.4.

\begin{table}[h]
\centering
\footnotesize
\caption{Human-label cross-rater audit under source-speech-group-held-out folds. Inputs are the stated rater's LQ/EXP plus segment-onset delay; brackets are speech-group-clustered 95\% CIs.}
\label{tab:cross-rater}
\begin{tabular}{lcc}
\toprule
Quality inputs & Promptness target & Pearson $r$ \\
\midrule
R05 & R05 & \CrossRtoR~[\CrossRtoRCILow, \CrossRtoRCIHigh] \\
R05 & R06 & \CrossRtoS~[\CrossRtoSCILow, \CrossRtoSCIHigh] \\
R06 & R05 & \CrossStoR~[\CrossStoRCILow, \CrossStoRCIHigh] \\
R06 & R06 & \CrossStoS~[\CrossStoSCILow, \CrossStoSCIHigh] \\
\bottomrule
\end{tabular}
\end{table}

\paragraph{Interpreter-disjoint results.} Table~\ref{tab:loio} reports the
pooled and macro-interpreter results summarized in Section~5.5.

\begin{table}[h]
\centering
\footnotesize
\caption{Interpreter-disjoint prediction of the aggregate target. Automatic rows report mean $\pm$ seed SD; macro $r$ equally weights held-out interpreters.}
\label{tab:loio}
\begin{tabular}{lrr}
\toprule
System & Pooled $r$ & Macro $r$ \\
\midrule
Delay & \LOIODelayPearson & \LOIODelayMacroPearson \\
Pred. Q + delay & $\LOIOQualityDelayPearson\pm\LOIOQualityDelayPearsonSD$ & $\LOIOQualityDelayMacroPearson\pm\LOIOQualityDelayMacroPearsonSD$ \\
Structure + delay & \LOIOStructureDelayPearson & \LOIOStructureDelayMacroPearson \\
Full model & $\mathbf{\LOIOFullPearson\pm\LOIOFullPearsonSD}$ & $\mathbf{\LOIOFullMacroPearson\pm\LOIOFullMacroPearsonSD}$ \\
\bottomrule
\end{tabular}
\end{table}

\paragraph{Direction-specific automatic systems.} Table~\ref{tab:automatic-direction}
reports the automatic-system sensitivity results summarized in Section~5.6.

% Generated by build_eacl_paper_results.py. Do not edit by hand.
\begin{table*}[t]
\centering
\scriptsize
\caption{Direction-specific results for the actual automatic systems. Seeded rows report mean $\pm$ sample SD across three fixed training seeds; deterministic rows are shown once. These descriptive subgroups are imbalanced and are not direction-level significance tests.}
\label{tab:automatic-direction}
\begin{tabular}{llccc}
\toprule
System & Direction ($n$) & Pearson $r$ & Spearman $r_S$ & MSE \\
\midrule
Segment-onset delay & Zh$\rightarrow$En (505) & .337 & .364 & .243 \\
Segment-onset delay & En$\rightarrow$Zh (117) & .319 & .305 & .446 \\
Predicted LQ+EXP + delay & Zh$\rightarrow$En (505) & $.414\pm.022$ & $.452\pm.013$ & $.235\pm.007$ \\
Predicted LQ+EXP + delay & En$\rightarrow$Zh (117) & $.717\pm.037$ & $.786\pm.012$ & $.295\pm.021$ \\
Lexical/structural + delay & Zh$\rightarrow$En (505) & .575 & .557 & .181 \\
Lexical/structural + delay & En$\rightarrow$Zh (117) & .749 & .790 & .211 \\
Full model & Zh$\rightarrow$En (505) & $.600\pm.013$ & $.592\pm.007$ & $.174\pm.004$ \\
Full model & En$\rightarrow$Zh (117) & $.777\pm.010$ & $.820\pm.006$ & $.192\pm.006$ \\
\bottomrule
\end{tabular}
\end{table*}


\paragraph{Structural residualization.} For each source-speech-group-held-out
fold and seed, a structure-to-predicted-quality mapping was fit on outer-training
rows only, then applied to form train/test quality residuals. Structural features
explain only a limited proportion of variation in the predicted quality outputs.
Because raw and residualized quality span closely related feature spaces when
structural features are retained, near-identical full-model predictions follow
from linear feature-space equivalence. Table~\ref{tab:residualization}
adds two less algebraically redundant checks: residualized-quality-only prediction,
and prediction of the structure+delay promptness residual from quality residuals.
The latter gives a small held-out partial $R^2$, while every seed-specific paired
$\Delta r$ interval crosses zero.

% Generated by build_eacl_paper_results.py. Do not edit by hand.
\begin{table*}[t]
\centering
\scriptsize
\caption{Fold-safe residualization representation audit. $R^2_Q$ gives structure-to-predicted-LQ/EXP held-out $R^2$; model cells report Pearson $r$/MSE. Partial $R^2$ is the held-out error reduction when quality residuals predict the structure+delay professional-promptness residual.}
\label{tab:residualization}
\begin{tabular}{lcccc}
\toprule
Seed & $R^2_Q$ LQ/EXP & Residual Q only & Residual augmentation & Partial $R^2$ \\
\midrule
20260718 & .200/.199 & .197/.305 & .664/.176 & .057 \\
20260719 & .246/.114 & .091/.319 & .659/.176 & .058 \\
20260720 & .215/.144 & .113/.315 & .645/.183 & .020 \\
\bottomrule
\end{tabular}
\end{table*}


\paragraph{Cross-rater result provenance.} The values in Table~\ref{tab:cross-rater}
use a literal target-specific specification: for each ordered pair, the stated
quality rater's LQ/EXP plus delay are used to fit the stated target rater's
promptness scores within the outer-training speech groups. Earlier working values
of $.502$ and $.828$ came from a different domain-swap analysis that fitted each
rater to that same rater's promptness and substituted the other rater's feature
scale and target only at held-out evaluation. They are therefore superseded and
are not comparable estimates of the cross-rater matrix reported here.

\paragraph{Direction-specific rater audit.} Table~\ref{tab:rater-direction}
reports the two most informative human-label systems by direction. The
English-to-Chinese condition contains only 117 segments, so these intervals are
descriptive rather than a direction-level model comparison.

% Generated by build_eacl_paper_results.py. Do not edit by hand.
\begin{table*}[t]
\centering
\scriptsize
\caption{Direction-specific human-label sensitivity. Brackets are speech-group clustered 95\% CIs.}
\label{tab:rater-direction}
\begin{tabular}{llrrr}
\toprule
Target & System & Direction ($n$) & Pearson $r$ & MSE \\
\midrule
R05 & LQ+EXP+delay & Zh$\rightarrow$En (505) & .895 [.814,.948] & .103 \\
R05 & LQ+EXP+delay & En$\rightarrow$Zh (117) & .929 [.860,.981] & .157 \\
R05 & Structure+delay & Zh$\rightarrow$En (505) & .508 [.262,.585] & .384 \\
R05 & Structure+delay & En$\rightarrow$Zh (117) & .740 [.663,.775] & .503 \\
R06 & LQ+EXP+delay & Zh$\rightarrow$En (505) & .537 [.466,.601] & .258 \\
R06 & LQ+EXP+delay & En$\rightarrow$Zh (117) & .641 [.464,.732] & .173 \\
R06 & Structure+delay & Zh$\rightarrow$En (505) & .428 [.342,.486] & .295 \\
R06 & Structure+delay & En$\rightarrow$Zh (117) & .466 [.372,.548] & .221 \\
\bottomrule
\end{tabular}
\end{table*}


\paragraph{Reproducibility inventory.} The camera-ready package will include the
corrected-data release tag and correction log, a table-to-file provenance map,
the exact outer and inner fold manifests, train/development/test counts,
environment versions, model-card configuration, calibration/range audit,
direction and per-interpreter tables, and the qualitative case-selection rule.
The present anonymous review version omits identifying links and restricted raw
media or identity mappings. The numerical source for the manuscript is the
canonical \texttt{paper\_results.json}, generated by
\texttt{build\_eacl\_paper\_results.py}; the cross-rater estimand correction is
documented in the reproducibility correction log.

\section*{Appendix C: Illustrative Cases}

Table~\ref{tab:case-pairs} gives one positive and one negative observed case.
They were selected post hoc under the documented purposive rule in the
documented case-selection rule; comments are translated and shortened, and are not
model inputs. They illustrate how evaluators articulated their ratings, but do not independently validate the promptness construct.

\begin{center}
\begin{minipage}{\columnwidth}
\centering
\scriptsize
\renewcommand{\arraystretch}{1.12}
\setlength{\tabcolsep}{2pt}
\captionof{table}{Observed professional-comment examples. These are illustrations, not additional rubric anchors or quantitative evidence.}
\label{tab:case-pairs}
\begin{tabular}{p{.14\columnwidth}p{.76\columnwidth}}
\toprule
Case & Observed example \\
\midrule
Positive & \textbf{Source:} An English conditional explains that many children in orphanages are not orphans, so ``orphanage'' is a euphemistic label. \newline \textbf{Interpreted output:} The Chinese output retains both the condition and conclusion. \newline \textbf{Scores:} Promptness $3.0$; predicted promptness $3.05$; segment-onset delay 2.00~s. \newline \textbf{Comment:} ``Fast, steady, complete, and idiomatic.'' \\
Negative & \textbf{Source:} An English segment states that a Mars trip requires an Earth--Mars alignment launch window. \newline \textbf{Interpreted output:} The Chinese output gives an incorrect distance and omits the launch-window condition. \newline \textbf{Scores:} Promptness $0.5$; predicted promptness $2.35$; segment-onset delay 2.44~s. \newline \textbf{Comment:} ``Limited information capture; unable to form complete sentences.'' \\
\bottomrule
\end{tabular}
\renewcommand{\arraystretch}{1}
\setlength{\tabcolsep}{6pt}
\end{minipage}
\end{center}

\section*{Appendix D: Quality Estimator Model Card}

The upstream estimator maps a complete ordered pair of source and
interpreted-output transcripts to two continuous quality outputs. It does not
use audio, delay, professional scores, rater/interpreter identity, or a
reference translation, and it is a post-hoc segment model rather than a
streaming system.

\begin{center}
\begin{minipage}{\columnwidth}
\centering
\scriptsize
\setlength{\tabcolsep}{2pt}
\captionof{table}{Configuration of the cross-fitted quality estimator used in the reported runs.}
\label{tab:model-card}
\begin{tabular}{p{.28\columnwidth}p{.63\columnwidth}}
\toprule
Item & Formal configuration \\
\midrule
Backbone & \texttt{Unbabel/wmt22-cometkiwi-da}; \texttt{unbabel-comet==2.2.4} loader \\
Input & Source as sequence A, interpreted output as sequence B; maximum 512 tokens; max-length padding and truncation \\
Tokenizer & Checkpoint tokenizer when available; otherwise \texttt{microsoft/infoxlm-large} \texttt{XLMRobertaTokenizer}. Logs do not preserve which branch or the checkpoint revision hash. \\
Representation & First-token (CLS) pooling; two independent linear residual heads added to fold-training LQ/EXP means \\
Trainable parameters & Two heads only. The instantiated LoRA adapters (rank 8, $\alpha=16$, dropout .1, query/value targets) and backbone remain frozen. \\
Optimization & AdamW; head LR $5\times10^{-4}$; weight decay .01; batch 16; 10 fixed epochs; no scheduler or early stopping; gradient norm .5; CUDA mixed precision when available \\
Selection and seeds & Maximum development Pearson(LQ)+Pearson(EXP); seeds 20260718, 20260719, and 20260720 \\
\bottomrule
\end{tabular}
\setlength{\tabcolsep}{6pt}
\end{minipage}
\end{center}

For batch predictions $\hat q_L,\hat q_E$ and targets $q_L,q_E$, training
minimizes
\begin{equation}
\begin{aligned}
\mathcal L={}&\mathrm{MSE}(\hat q_L,q_L)+1.7\,\mathrm{MSE}(\hat q_E,q_E)\\
&+0.2\,\mathcal L_{\mathrm{var}}\\
&+0.25[1-r(\hat q_L,q_L)]\\
&+0.25[1-r(\hat q_E,q_E)].
\end{aligned}
\end{equation}
The variance term matches prediction variance to target variance using a
rolling 64-example buffer and a .15 target-variance floor. These anti-collapse
weights were fixed across reported runs; the archive does not preserve a
separate search that selected them. Development labels are rounded to the
nearest .5 only in the historical validation routine, whose predictions are
clamped to $[-1,4]$ for checkpoint-selection correlations; raw exported quality
predictions are unchanged. Transcript production/editing provenance,
historical truncation counts, exact hardware, and per-fold training time are
not preserved.
